\uuid{MklJ}
\titre{Raisonnement}
\auteur{bodin ; hanani ; mzari}
\organisation{exo7}
\licence{CC BY-NC-SA 4.0}
\datecreate{2020-01-20}
\niveau{L1}
\module{Algèbre}
\chapitre{Logique, ensemble, raisonnement}
\sousChapitre{Absurde et contraposée}
\difficulte{1}
\theme{Raisonnement}
\exo7qcmid{lille-1.logique.019}

\contenu{
\begin{question}
\qclassification{100.03,100.04}
\qtags{motcle=raisonnement}
Je veux montrer que $\frac{n(n+1)}{2}$ est un entier, quelque soit $n\in\Nn$.  Quelles sont les démarches possibles ?
\begin{answers}
    \bad{Montrer que la fonction $x \mapsto x(x+1)$ est paire.}

    \good{Séparer le cas $n$ pair, du cas $n$ impair.}

    \bad{Par l'absurde, supposer que $\frac{n(n+1)}{2}$ est un réel, puis chercher une contradiction.}

    \bad{Le résultat est faux, je cherche un contre-exemple.}
\end{answers}
\begin{explanations}
Séparer le cas $n$ pair, du cas $n$ impair. Dans le premier cas, on peut écrire $n=2k$ (avec $k\in \Nn$), dans le second cas $n=2k+1$, puis calculer $\frac{n(n+1)}{2}$.
\end{explanations}
\end{question}
}
